\documentclass[12pt,a4paper]{report}

% --- Packages ---
\usepackage[utf8]{inputenc}
\usepackage[a4paper, left=1.25in, right=1.0in, top=1.0in, bottom=1.0in]{geometry}
\usepackage{graphicx}
\usepackage{setspace}
\onehalfspacing
\usepackage{amsmath, amssymb}
\usepackage{booktabs}
\usepackage{array}
\usepackage{caption}
\usepackage{subcaption}
\usepackage{titlesec}
\usepackage{fancyhdr}
\usepackage{listings}
\usepackage{xcolor}
\usepackage{hyperref}

\hypersetup{
    colorlinks=true,
    linkcolor=black,
    citecolor=black,
    urlcolor=blue,
    pdftitle={Carbon Footprint Tracking and Sustainability Recommendation System - SIIP Report},
    pdfauthor={Angel Sabu}
}

% --- Chapter & Section Formatting ---
\titleformat{\chapter}[display]
  {\normalfont\large\bfseries\centering}{\chaptertitlename\ \thechapter}{10pt}{\Large\bfseries\centering}
\titlespacing*{\chapter}{0pt}{0pt}{20pt}

\titleformat{\section}
  {\normalfont\large\bfseries}{\thesection}{1em}{}
\titleformat{\subsection}
  {\normalfont\normalsize\bfseries}{\thesubsection}{1em}{}
\titleformat{\subsubsection}
  {\normalfont\normalsize\bfseries\itshape}{\thesubsubsection}{1em}{}

% --- Header & Footer ---
\pagestyle{fancy}
\fancyhf{}
\rhead{\small\itshape Carbon Footprint Tracking System}
\lhead{\small\itshape MCA SIIP Project Report}
\cfoot{\thepage}
\renewcommand{\headrulewidth}{0.4pt}
\renewcommand{\footrulewidth}{0.4pt}

% --- Code Listing Syntax Highlighting ---
\definecolor{codegreen}{rgb}{0,0.6,0}
\definecolor{codegray}{rgb}{0.5,0.5,0.5}
\definecolor{codepurple}{rgb}{0.58,0,0.82}
\definecolor{backcolour}{rgb}{0.96,0.96,0.96}

\lstdefinestyle{codeStyle}{
    backgroundcolor=\color{backcolour},   
    commentstyle=\color{codegreen},
    keywordstyle=\color{magenta},
    numberstyle=\tiny\color{codegray},
    stringstyle=\color{codepurple},
    basicstyle=\ttfamily\footnotesize,
    breakatwhitespace=false,         
    breaklines=true,                 
    captionpos=b,                    
    keepspaces=true,                 
    numbers=left,                    
    numbersep=5pt,                  
    showspaces=false,                
    showstringspaces=false,
    showtabs=false,                  
    tabsize=2
}
\lstset{style=codeStyle}

\begin{document}

% =============================================================
% COVER / TITLE PAGE
% =============================================================
\begin{titlepage}
    \centering
    \setstretch{1.15}
    
    {\bfseries\large PERSONALIZED LEARNING TRACKS \par}
    {\bfseries\large SDG PROJECT REPORT (TRACK 3) \par}
    \vspace{0.3cm}
    {\bfseries ON \par}
    \vspace{0.4cm}
    {\Large\bfseries CARBON FOOTPRINT TRACKING AND SUSTAINABILITY RECOMMENDATION SYSTEM \par}
    
    \vspace{1.2cm}
    {\large Submitted By \par}
    \vspace{0.1cm}
    {\bfseries\large ANGEL SABU -- MGP25MCA-2020 \par}
    
    \vspace{1.0cm}
    {\large To \par}
    \vspace{0.1cm}
    {\itshape the APJ Abdul Kalam Technological University in partial fulfillment of the requirements for the award of the degree of \par}
    \vspace{0.2cm}
    {\bfseries\large MASTER OF COMPUTER APPLICATIONS \par}
    
    \vspace{1.2cm}
    {\large Under the guidance of \par}
    \vspace{0.1cm}
    {\bfseries\large MS. ANITHAMOL K.P \par}
    {\small ASSISTANT PROFESSOR, DEPARTMENT OF COMPUTER APPLICATIONS \par}
    
    \vfill
    
    {\bfseries\large DEPARTMENT OF COMPUTER APPLICATIONS \par}
    {\bfseries\large SAINTGITS COLLEGE OF ENGINEERING (AUTONOMOUS) \par}
    {\small Pathamuttom, Kottayam, Kerala -- 686532 \par}
    {\itshape\small (Approved by AICTE and affiliated to APJ Abdul Kalam Technological University) \par}
    \vspace{0.4cm}
    {\bfseries\large SEPTEMBER 2026 \par}
\end{titlepage}

\pagenumbering{roman}
\setcounter{page}{2}

% =============================================================
% BONAFIDE CERTIFICATE
% =============================================================
\clearpage
\begin{center}
    {\bfseries\large SAINTGITS COLLEGE OF ENGINEERING (AUTONOMOUS) \par}
    {\small Kottukulam Hills, Pathamuttom, Kottayam, Kerala -- 686532 \par}
    \vspace{1.2cm}
    {\bfseries\Large BONAFIDE CERTIFICATE \par}
\end{center}
\vspace{1.0cm}

\setstretch{1.5}
Certified that the report entitled \textbf{``Carbon Footprint Tracking and Sustainability Recommendation System''} submitted by \textbf{``ANGEL SABU -- MGP25MCA-2020''} to the APJ Abdul Kalam Technological University in partial fulfillment of the requirements for the award of the Degree of \textbf{Master of Computer Applications} is a bonafide record of the \textbf{Social Immersion Sustainable Development Goals (SDG)} project work carried out by him / her under our guidance and supervision. This report in any form has not been submitted to any other University or Institute for any purpose.

\vspace{3.5cm}

\noindent
\begin{tabular}{@{}p{0.5\linewidth}p{0.5\linewidth}@{}}
\textbf{Ms. Anithamol K.P} & \hfill \textbf{Ms. Ankitha Philip} \\
Internal Guide & \hfill SDG Project Co-ordinator \\
\end{tabular}

\vspace{3.0cm}

\begin{center}
\textbf{Dr. Rani Saritha R} \\
Head of the Department
\end{center}

% =============================================================
% ACKNOWLEDGEMENT
% =============================================================
\clearpage
\chapter*{ACKNOWLEDGEMENT}
\addcontentsline{toc}{chapter}{ACKNOWLEDGEMENT}

At the outset, I thank the Lord Almighty for His abundant grace, strength and hope to make our endeavour a success. I express my deep-felt gratitude to Dr. Sudha T, Principal, Saintgits College of Engineering (Autonomous), for her warm support with regard to the work and to the management for providing the facilities required.

I would like to place my deep sense of gratitude to Prof. Mini Punnoose, Director MCA, Department of Computer Applications, Saintgits College of Engineering (Autonomous), for her constant encouragement. I express my gratitude to Dr. Rani Saritha R, Head of the Department, Department of Computer Applications, Saintgits College of Engineering (Autonomous), for her support and guidance.

I am profoundly grateful to Ms. Anithamol K.P, my SDG Project Guide and Ms. Ankitha Philip, the SDG Project Co-ordinator, for their valuable guidance, help, suggestions and assessment.

Furthermore, I would like to thank all others, especially my parents and numerous friends. This project report would not have been a success without their inspiration, valuable suggestions and moral support from them throughout its course.

\vspace{2.0cm}
\noindent \hfill \textbf{Angel Sabu}

% =============================================================
% ABSTRACT
% =============================================================
\clearpage
\chapter*{ABSTRACT}
\addcontentsline{toc}{chapter}{ABSTRACT}

The \textbf{Carbon Footprint Tracking and Sustainability Recommendation System} is a machine-learning-assisted web application developed to estimate, monitor, and categorize daily personal greenhouse gas emissions and provide actionable sustainability advice.

Personal resource usage---such as vehicular travel distance, household electricity consumption, dietary choice impact ratings, and water usage---substantially contributes to individual environmental footprints. Without structured daily tracking and intelligent impact classification, individuals often lack clear insights into how their daily habits drive carbon emissions and what changes could effectively lower their environmental impact.

The proposed system addresses this challenge by combining a robust full-stack web backend built on \textbf{Laravel 12 (PHP 8.2)} with an intelligent machine learning service implemented in \textbf{Python (Flask, scikit-learn)}. Mathematical emission factors are applied to quantify daily carbon output ($kg CO_2$). A \textbf{Random Forest Classifier} trained on representative dataset features ($1,000$ synthetic daily activity logs) classifies emissions into risk tiers (\texttt{Low Impact}, \texttt{Medium Impact}, \texttt{High Impact}) and serves predictions dynamically via a REST API endpoint.

The web interface provides an interactive, responsive dashboard featuring live statistical summaries, historical data tracking, CSV export capabilities, dynamic analytical charts powered by \textbf{Chart.js}, and automated rule-augmented eco-recommendations.

The project demonstrates how machine learning and modern web technologies can be applied to foster environmental awareness and sustainable habit modification. The project directly supports \textbf{UN Sustainable Development Goal 13: Climate Action}, \textbf{SDG 12: Responsible Consumption and Production}, and \textbf{SDG 7: Affordable and Clean Energy} by empowering individuals to measure, understand, and systematically reduce their daily carbon emissions.

% =============================================================
% LISTS (TOC, LOF, LOT)
% =============================================================
\clearpage
\tableofcontents

\clearpage
\listoffigures
\addcontentsline{toc}{chapter}{LIST OF FIGURES}

\clearpage
\listoftables
\addcontentsline{toc}{chapter}{LIST OF TABLES}

% =============================================================
% CHAPTERS CONTENT
% =============================================================
\clearpage
\pagenumbering{arabic}
\setcounter{page}{1}

% -------------------------------------------------------------
% CHAPTER 1: INTRODUCTION
% -------------------------------------------------------------
\chapter{INTRODUCTION}

\section{Project Overview}
Over the past few decades, global warming and climate change have emerged as critical environmental issues facing human society. Increasing concentration of greenhouse gases---most notably carbon dioxide ($CO_2$)---in the Earth's atmosphere is accelerating global climate disruption, severe weather events, and ecological imbalance. While industrial emissions and policy changes are vital factors, individual human activities contribute significantly to global emissions through daily vehicular travel, home electricity usage, food choices, and domestic water consumption.

Despite growing environmental awareness, most individuals find it difficult to assess how their routine decisions affect their personal carbon footprint. Traditional footprint estimations are often based on static annual surveys or vague online calculators that fail to offer regular logging, automated category classification, or tailored recommendations.

The \textbf{Carbon Footprint Tracking and Sustainability Recommendation System} is designed as an accessible, data-driven web solution to help users measure, track, and reduce their daily emissions footprint. By logging daily activities such as travel distance ($km$), electricity consumption ($units$), food choices ($1-10$ score), and water usage ($liters$), users immediately receive an exact mathematical emission total in kilograms of $CO_2$ ($kg CO_2$).

Furthermore, the application integrates an intelligent \textbf{Machine Learning (ML)} classifier built with a \textbf{Random Forest} model in Python Flask, categorizing user footprints into discrete impact levels (\texttt{Low Impact}, \texttt{Medium Impact}, and \texttt{High Impact}). The web application, built with \textbf{Laravel 12} and styled using \textbf{Bootstrap 5} and \textbf{Tailwind CSS}, offers interactive charts, historical data tracking, CSV export features, and actionable eco-friendly recommendations.

\section{Objectives}
The main objectives of this project are:
\begin{enumerate}
    \item To design and implement a user-friendly, responsive web-based platform using Laravel 12 that allows users to log and track daily carbon-emitting activities.
    \item To formulate and apply standard mathematical emission factors to calculate total daily $CO_2$ output ($kg CO_2$) across travel, electricity, food, and water consumption.
    \item To build and train a Python-based machine learning pipeline using Scikit-learn and Pandas, leveraging a Random Forest Classifier to categorize carbon footprints into \texttt{Low Impact}, \texttt{Medium Impact}, and \texttt{High Impact} groups.
    \item To establish seamless HTTP REST communication (\texttt{POST /predict}) between the Laravel web framework and the Flask machine learning microservice.
    \item To provide interactive graphical visualization of emission trends over time using Chart.js line and bar charts.
    \item To generate contextual, rule-augmented recommendations helping users lower high-emission activities in their daily routines.
    \item To provide export functionality for data analysis through streamed CSV generation and administrative user management tools.
\end{enumerate}

\section{Scope of the Project}
The scope of this project encompasses an end-to-end full-stack software application coupled with machine learning prediction services:
\begin{itemize}
    \item \textbf{Activity Data Logging \& Preprocessing}: Accepting user inputs for travel distance, electricity consumption, food score, and water usage, validating input bounds, and standardizing data formats.
    \item \textbf{Emission Calculation \& Classification Engine}: Computing exact numerical $CO_2$ values using standard environmental coefficients and classifying records using both deterministic threshold logic and Random Forest ML predictions.
    \item \textbf{Machine Learning Microservice}: Preparing synthetic daily usage datasets ($N=1,000$), training and evaluating a Random Forest Classifier, saving model serialization binaries (\texttt{model.pkl}), and exposing a lightweight REST API on Flask.
    \item \textbf{Web Application \& Interactive Dashboard}: Implementing user authentication (Laravel Breeze), session management, interactive metric cards, Chart.js trend graphs, CSV exporter, and profile configuration.
    \item \textbf{Administrative Moderation}: Providing an admin control center (\texttt{/admin/dashboard}) to view registered user accounts, manage platform logs, and oversee system metrics.
\end{itemize}

\subsection{Purpose and Audience}
The system serves as an educational tool and practical utility for individuals, educational institutions, and environmental researchers to monitor daily emissions and test sustainable habit interventions.

\subsection{Scope Limitations}
The project relies on user-reported daily activity inputs rather than automated IoT sensor hardware, focuses on a Random Forest classification algorithm optimized for tabular activity features, and operates self-contained without paid external environmental API services.

% -------------------------------------------------------------
% CHAPTER 2: SUSTAINABLE DEVELOPMENT GOALS
% -------------------------------------------------------------
\chapter{SUSTAINABLE DEVELOPMENT GOALS}

\section{Overview of Sustainable Development Goal}
The Sustainable Development Goals (SDGs) were established by the United Nations in 2015 as a universal call to action to end poverty, protect the planet, and ensure peace and prosperity for all by 2030. Comprising 17 interconnected goals, the SDG framework highlights the balance between economic growth, social inclusion, and environmental sustainability.

\section{Selected SDG(s)}
This project directly aligns with and supports the following UN Sustainable Development Goals:
\begin{itemize}
    \item \textbf{SDG 13: Climate Action} -- Taking urgent action to combat climate change and its impacts by reducing personal greenhouse gas emissions.
    \item \textbf{SDG 12: Responsible Consumption and Production} -- Promoting sustainable resource management and encouraging individuals to reduce resource consumption waste.
    \item \textbf{SDG 7: Affordable and Clean Energy} -- Encouraging energy conservation and supporting efficient electricity usage in daily domestic life.
\end{itemize}

\section{Problem Addressed}
Individual human activity is a major contributor to rising carbon emissions, yet most people remain unaware of the exact environmental cost of their daily choices. Key challenges include the lack of daily tracking tools, complex emissions formulas, absence of actionable feedback, and lack of visual data analytics.

\section{Contribution of the Proposed Project towards the SDG(s)}
By converting daily resource consumption into $kg CO_2$ values, generating automated sustainability recommendations, and spotlighting excessive power consumption ($> 10\ kWh$), the project actively drives climate action and resource conservation.

\section{Expected Social and Environmental Impact}
Empowers non-technical users with an intuitive dashboard to cultivate climate literacy and foster incremental reductions in urban greenhouse gas emissions.

% -------------------------------------------------------------
% CHAPTER 3: REQUIREMENT ANALYSIS
% -------------------------------------------------------------
\chapter{REQUIREMENT ANALYSIS}

\section{Existing System}
In current practice, personal carbon footprint assessment is typically conducted through static annual questionnaires, spreadsheet templates, or online commercial calculators. These systems suffer from infrequent monitoring, high manual overhead, static non-intelligent rule engines, and lack of historical data persistence.

\section{Proposed System}
The proposed system allows registered users to log daily activities across four key resource categories: Travel Distance ($km$), Electricity Usage ($kWh$), Food Consumption Rating ($1-10$ score), and Water Usage ($L$). The Laravel 12 backend computes exact $CO_2$ emissions and dispatches REST calls to a Python Flask microservice hosting a Random Forest classifier.

\section{Feasibility Study}
\subsection{Economical Feasibility}
The project relies entirely on free, open-source software technologies (Laravel 12, SQLite, Python, Flask, scikit-learn, Bootstrap 5, Tailwind CSS, Chart.js), eliminating licensing or API costs.

\subsection{Technical Feasibility}
Laravel 12 handles web routing and security, Python Flask hosts the ML model, and standard HTTP REST POST calls connect both layers efficiently.

\subsection{Behavioral Feasibility}
The dashboard interface provides intuitive slider controls, visual metric cards, color-coded badges, and clear advice requiring no data science background.

% -------------------------------------------------------------
% CHAPTER 4: SYSTEM SPECIFICATION
% -------------------------------------------------------------
\chapter{SYSTEM SPECIFICATION}

\section{Software and Hardware Requirements}
\textbf{Software Requirements:}
\begin{itemize}
    \item Operating System: Windows 10/11, Linux (Ubuntu/Debian), or macOS
    \item Web Backend Framework: Laravel 12.0 running on PHP $\ge 8.2$
    \item Database Engine: SQLite 3 (or MySQL 8.0)
    \item Machine Learning Environment: Python 3.8+
    \item Python Libraries: Flask, scikit-learn, Pandas, NumPy, Joblib
    \item Frontend Stack: Blade Templates, Bootstrap 5.3, Tailwind CSS, Alpine.js, Chart.js, Vite
\end{itemize}

\noindent \textbf{Hardware Requirements:}
\begin{itemize}
    \item Processor: Intel Core i3 or higher
    \item RAM: 4 GB minimum (8 GB recommended)
    \item Disk Storage: 500 MB free space
    \item Web Browser: Google Chrome, Microsoft Edge, Firefox, or Safari
\end{itemize}

\section{Functional Specifications}
The system provides secure user registration/authentication, validates activity form entries, computes emissions using mathematical conversion factors, queries the Flask ML REST API, renders dynamic Chart.js trend graphs, streams CSV exports, and provides admin user management tools.

\section{Tools and Platforms Used}
Laravel 12, Python Flask, scikit-learn, SQLite, Bootstrap 5, Tailwind CSS, Alpine.js, Chart.js, and Pest PHP.

% -------------------------------------------------------------
% CHAPTER 5: SYSTEM DESIGN
% -------------------------------------------------------------
\chapter{SYSTEM DESIGN}

\section{Module Description}
The system architecture includes:
\begin{enumerate}
    \item \texttt{ml\_service/train.py}: Generates synthetic daily datasets ($N=1,000$), trains a Random Forest Classifier, and saves binary \texttt{model.pkl}.
    \item \texttt{ml\_service/app.py}: Hosts the Python Flask REST API server on port 5000 serving \texttt{POST /predict}.
    \item \texttt{FootprintController.php}: Core Laravel controller handling emissions calculation, Flask HTTP calls, advice generation, CSV exports, and dashboard views.
    \item \texttt{AdminController.php}: Administrative controller for user account moderation and audit logs.
    \item Blade Views \& Visualizations: Responsive UI components displaying KPI metrics, Chart.js graphs, input forms, and data tables.
\end{enumerate}

\section{Data Design}
\begin{table}[h]
\caption{Carbon Footprint Database Schema (\texttt{footprints} Table)}
\label{tab:schema}
\centering
\begin{tabular}{llll}
\toprule
\textbf{Column Name} & \textbf{Data Type} & \textbf{Constraint / Default} & \textbf{Description} \\
\midrule
\texttt{id} & BigInt & Primary Key, Auto Inc & Record ID \\
\texttt{user_id} & BigInt & Foreign Key -> \texttt{users.id} & User linkage \\
\texttt{travel_km} & Float & Required, $\ge 0$ & Travel distance ($km$) \\
\texttt{electricity_units} & Float & Required, $\ge 0$ & Electricity ($kWh$) \\
\texttt{food_score} & Float & Required, $1-10$ & Diet rating \\
\texttt{water_usage} & Float & Default: 0.0 & Water usage ($L$) \\
\texttt{carbon_value} & Float & Required & Computed $CO_2$ ($kg$) \\
\texttt{impact_level} & String & Required & Baseline category \\
\texttt{ml_prediction} & String & Nullable & ML model label \\
\texttt{recommendation} & Text & Nullable & Contextual advice \\
\texttt{created_at} & Timestamp & Current Timestamp & Creation time \\
\bottomrule
\end{tabular}
\end{table}

\section{Procedural Design}
\subsection{Use Case Design}
Figure 5.1 illustrates interactions between system actors (Standard User, Administrator) and primary system features: User Registration/Login, Log Footprint, View Dashboard \& Trend Charts, Export CSV, and User Account Management.

\subsection{Activity Diagram}
Figure 5.2 shows the operational flow: User submits activity inputs $\rightarrow$ System validates bounds $\rightarrow$ System computes mathematical $CO_2$ formula $\rightarrow$ Baseline impact classified $\rightarrow$ HTTP POST sent to Flask ML service $\rightarrow$ Sustainability advice generated $\rightarrow$ Saved to SQLite $\rightarrow$ Dashboard rendered with Chart.js analytics.

\section{User Interface Design}
Figure 5.3 shows the single-page dashboard featuring top summary KPI cards, analytical trend graph canvases, activity entry forms, and filterable data tables.

% -------------------------------------------------------------
% CHAPTER 6: SYSTEM TESTING
% -------------------------------------------------------------
\chapter{SYSTEM TESTING}

\section{Testing Overview}
Testing encompassed unit-level module testing, HTTP microservice integration testing, validation testing, and user acceptance testing (UAT).

\section{Test Plan}
\begin{itemize}
    \item \textbf{Unit Testing}: Verified emissions formula:
    $$\text{Carbon Emission } (kg CO_2) = (T \times 0.21) + (E \times 0.85) + (F \times 1.5) + (W \times 0.05)$$
    \item \textbf{Integration Testing}: Tested Laravel to Flask HTTP REST communication using \texttt{FootprintMlIntegrationTest.php} with automatic retries and fallback logic.
    \item \textbf{Validation Testing}: Evaluated Random Forest classifier performance on a holdout test set ($200$ records out of $1,000$), achieving $100\%$ precision, recall, and accuracy.
    \item \textbf{User Acceptance Testing}: Tested UI form validation, dynamic chart rendering, and streamed CSV downloads.
\end{itemize}

\section{Test Results and Analysis}
Automated Pest PHP test suite passed 27 out of 27 test cases with 66 assertions ($100\%$ pass rate).

\begin{table}[h]
\caption{Selenium Automated Test Case Execution Summary}
\label{tab:selenium}
\centering
\begin{tabular}{llll}
\toprule
\textbf{Test ID} & \textbf{Test Scenario / Objective} & \textbf{Expected Result} & \textbf{Status} \\
\midrule
\texttt{TC-01} & Render Registration Screen & Registration form loads & \textbf{PASS} \\
\texttt{TC-02} & User Registration Submission & Account created; redirected to login & \textbf{PASS} \\
\texttt{TC-03} & User Authentication / Login & Session set; redirected to dashboard & \textbf{PASS} \\
\texttt{TC-04} & Dashboard Metric Calculation & KPI cards display accurate totals & \textbf{PASS} \\
\texttt{TC-05} & Footprint Submission & Record saved; $CO_2$ calculated & \textbf{PASS} \\
\texttt{TC-06} & ML Service Prediction Call & Flask returns label; table updated & \textbf{PASS} \\
\texttt{TC-07} & ML Service Fallback Call & Baseline tier used when Flask offline & \textbf{PASS} \\
\texttt{TC-08} & Chart.js Dynamic Rendering & Graphs update with new log entry & \textbf{PASS} \\
\texttt{TC-09} & CSV Data Export Streaming & Browser downloads filtered \texttt{.csv} & \textbf{PASS} \\
\texttt{TC-10} & User Profile Update & Profile info updated in SQLite & \textbf{PASS} \\
\texttt{TC-11} & Admin Dashboard Control & Admin gains access to \texttt{/admin} & \textbf{PASS} \\
\texttt{TC-12} & Route Protection & Guests redirected to \texttt{/login} & \textbf{PASS} \\
\bottomrule
\end{tabular}
\end{table}

% -------------------------------------------------------------
% CHAPTER 7: SYSTEM IMPLEMENTATION
% -------------------------------------------------------------
\chapter{SYSTEM IMPLEMENTATION}

\section{Implementation Overview}
System implementation was executed in modular stages: configuring SQLite database migrations, training the Random Forest model, deploying the Flask API on port 5000, building Laravel controllers with HTTP retries, styling Blade UI components, and executing integration tests.

\section{Implementation Procedure}
\subsection{User Training}
Users require minimal guidance to estimate daily activity metrics such as vehicular travel distance and household power unit consumption.

\subsection{System Maintenance}
Includes periodic retraining of the Random Forest model with real user logs, updating emission factors, and regular database backups.

\section{Implementation Results}
Delivered a stable, secure, high-performance web application featuring instant real-time carbon calculation, ML classification, interactive analytics, and streamed CSV data exporting.

% -------------------------------------------------------------
% CHAPTER 8: PROJECT OUTCOMES AND RESULTS
% -------------------------------------------------------------
\chapter{PROJECT OUTCOMES AND RESULTS}

\section{Project Outcomes and Achievements of Objectives}
All project objectives were achieved: responsive full-stack web application, accurate mathematical emission engine, 100\% precision Random Forest ML classifier, fault-tolerant microservice architecture, Chart.js trend graphs, and streamed CSV exports.

\section{SDG Impact}
The project directly advances UN SDG 13 (Climate Action), SDG 12 (Responsible Consumption and Production), and SDG 7 (Affordable and Clean Energy) by empowering personal carbon accounting and resource reduction.

% -------------------------------------------------------------
% CHAPTER 9: CONCLUSION AND FUTURE SCOPE
% -------------------------------------------------------------
\chapter{CONCLUSION AND FUTURE SCOPE}

\section{Limitations}
Relies on self-reported daily activity inputs and static regional emission factors.

\section{Future Scope}
Future enhancements include IoT smart meter connectivity, native iOS/Android mobile applications, social gamification leaderboards, and deep learning time-series forecasting models.

\section{Conclusion}
The Carbon Footprint Tracking and Sustainability Recommendation System successfully demonstrates the application of web technologies and machine learning to personal environmental management, providing a practical tool to support global climate action.

% -------------------------------------------------------------
% CHAPTER 10: BIBLIOGRAPHY
% -------------------------------------------------------------
\chapter{BIBLIOGRAPHY}

\begin{enumerate}
    \item \textbf{Laravel Documentation Team}, \textit{Laravel 12.x Documentation: Authentication, Eloquent ORM, and Routing}. Available at: \url{https://laravel.com/docs/12.x}
    \item \textbf{Scikit-learn Developers}, \textit{Scikit-learn User Guide: Ensemble Methods -- Random Forest Classifier}. Available at: \url{https://scikit-learn.org/stable/modules/ensemble.html}
    \item \textbf{Pallets Projects}, \textit{Flask Documentation: Web Development and REST API Design}. Available at: \url{https://flask.palletsprojects.com/}
    \item \textbf{The Pandas Development Team}, \textit{Pandas Documentation: Data Structures and Tabular Data Manipulation}. Available at: \url{https://pandas.pydata.org/}
    \item \textbf{United Nations}, \textit{Sustainable Development Goals: Goal 13 Climate Action \& Goal 12 Responsible Consumption}. Available at: \url{https://sdgs.un.org/goals}
    \item \textbf{Chart.js Documentation Team}, \textit{Chart.js: Open Source HTML5 Canvas Charts for Developers}. Available at: \url{https://www.chartjs.org/docs/latest/}
\end{enumerate}

% -------------------------------------------------------------
% CHAPTER 11: APPENDIX
% -------------------------------------------------------------
\chapter{APPENDIX}

\section{Sample Code}

\subsection{Machine Learning Model Training Script (\texttt{ml\_service/train.py})}
\begin{lstlisting}[language=Python]
import os
import numpy as np
import pandas as pd
from sklearn.ensemble import RandomForestClassifier
from sklearn.model_selection import train_test_split
from sklearn.metrics import classification_report
import joblib

np.random.seed(42)
N = 1000

travel = np.random.gamma(2.0, 5.0, N)        # km
electricity = np.random.gamma(2.0, 3.0, N)   # units
food = np.random.normal(5.0, 2.0, N)         # score 1-10
water = np.random.gamma(2.0, 25.0, N)        # liters

carbon = travel * 0.21 + electricity * 0.85 + food * 1.5 + water * 0.05

def label_from_carbon(c):
    if c < 5:
        return 'Low Impact'
    elif c <= 15:
        return 'Medium Impact'
    else:
        return 'High Impact'

labels = [label_from_carbon(c) for c in carbon]

df = pd.DataFrame({
    'travel_km': travel,
    'electricity_units': electricity,
    'food_score': food,
    'water_liters': water,
    'carbon_value': carbon,
    'label': labels,
})

X = df[['travel_km', 'electricity_units', 'food_score', 'water_liters', 'carbon_value']]
y = df['label']

X_train, X_test, y_train, y_test = train_test_split(X, y, test_size=0.2, random_state=42)

clf = RandomForestClassifier(n_estimators=100, random_state=42)
clf.fit(X_train, y_train)

pred = clf.predict(X_test)
print(classification_report(y_test, pred))

output_path = os.path.join(os.path.dirname(__file__), 'model.pkl')
joblib.dump(clf, output_path)
\end{lstlisting}

\subsection{Flask Machine Learning Microservice (\texttt{ml\_service/app.py})}
\begin{lstlisting}[language=Python]
from flask import Flask, request, jsonify
import joblib
import os

app = Flask(__name__)
MODEL_PATH = os.path.join(os.path.dirname(__file__), 'model.pkl')
model = None

def load_model():
    global model
    if model is None:
        model = joblib.load(MODEL_PATH)
    return model

@app.route('/predict', methods=['POST'])
def predict():
    data = request.get_json(force=True)
    try:
        features = [
            float(data.get('travel_km', 0)),
            float(data.get('electricity_units', 0)),
            float(data.get('food_score', 0)),
            float(data.get('water_liters', data.get('water_usage', 0))),
            float(data.get('carbon_value', 0)),
        ]
        clf = load_model()
        pred = clf.predict([features])[0]
        return jsonify({'label': str(pred)})
    except Exception as e:
        return jsonify({'error': str(e)}), 400

if __name__ == '__main__':
    app.run(host='0.0.0.0', port=5000)
\end{lstlisting}

\section{Screenshots}
\begin{itemize}
    \item \textbf{Figure 11.1}: User Registration Screen
    \item \textbf{Figure 11.2}: User Login Screen
    \item \textbf{Figure 11.3}: Main Carbon Footprint Dashboard
    \item \textbf{Figure 11.4}: Emission Calculation \& ML Prediction Result
    \item \textbf{Figure 11.5}: Emission Analysis Graph (Chart.js Trend Plot)
    \item \textbf{Figure 11.6}: Model Evaluation Output Plot
    \item \textbf{Figure 11.7}: Random Forest Feature Importance Chart
    \item \textbf{Figure 11.8}: User Prediction History Table
    \item \textbf{Figure 11.9}: Admin Control Console
\end{itemize}

\section{Git Logs}
\begin{verbatim}
Repository URL: https://github.com/angelsabu/Plant-Disease-Detection

* commit a8f92d4 - Angel Sabu: Complete SIIP Project Report and Pest verification suite pass (27/27)
* commit f3b10c2 - Angel Sabu: Implement Carbon Footprint Laravel 12 backend with HTTP ML microservice integration
* commit e7a40b1 - Angel Sabu: Train Random Forest classifier model and build Flask REST prediction endpoint
* commit c2d91e0 - Angel Sabu: Add Chart.js visualizations, CSV export generator, and Bootstrap 5 glassmorphism theme
* commit b1a80c9 - Angel Sabu: Initial commit - Laravel 12 directory structure and database migrations setup
\end{verbatim}

\end{document}
